% Informe final: Sistema Multi-Agente ReAct
% Compila con pdfLaTeX, XeLaTeX, LuaLaTeX o tectonic (ver informe/README.md).
% Cada sección está en secciones/: editar ahí el contenido y aquí solo la portada.
\documentclass[11pt,a4paper]{article}

\usepackage{iftex}
\ifPDFTeX
  \usepackage[utf8]{inputenc}
  \usepackage[T1]{fontenc}
  \usepackage{lmodern}
\else
  \usepackage{fontspec}
\fi
% es-noshorthands: sin él, babel vuelve activos < y > y rompe las flechas de TikZ.
% es-nolists: conserva las viñetas estándar en vez de las de babel.
\usepackage[spanish,es-noshorthands,es-nolists,es-tabla,es-nodecimaldot]{babel}
\usepackage[margin=2.5cm]{geometry}
\usepackage{microtype}
\usepackage{amsmath}
\usepackage{booktabs}
\usepackage{tabularx}
\usepackage{multirow}
\usepackage{array}
\usepackage[table]{xcolor}
\usepackage{enumitem}
\usepackage{caption}
\usepackage{fancyvrb}
\usepackage{tikz}
\usetikzlibrary{positioning,arrows.meta,shapes.geometric,babel}
\usepackage{pgfplots}
\pgfplotsset{compat=1.18}
\usepackage[hidelinks]{hyperref}

\setlist{itemsep=2pt,topsep=4pt}
\captionsetup{font=small,labelfont=bf}
\newcommand{\code}[1]{\texttt{#1}}
\definecolor{multi}{HTML}{2E6DB4}
\definecolor{naive}{HTML}{C8742B}

\title{%
  \textbf{Sistema Multi-Agente ReAct}\\[4pt]
  \large Auditoría de respuestas con agentes, herramientas y guardrails\\
  \large sobre WikiQA, GSM8K y MBPP}
% EDITAR: nombres, códigos, curso y docente.
\author{Integrante 1 \and Integrante 2 \and Integrante 3}
\date{Procesamiento de Lenguaje Natural --- Proyecto Final\\[2pt] Noviembre de 2026}

\begin{document}
\maketitle

\begin{abstract}
Presentamos un sistema multi-agente basado en el paradigma ReAct que coordina tres agentes en un grafo cíclico de LangGraph:
un \emph{Investigador}, que razona y usa herramientas con el ciclo \emph{Thought $\rightarrow$ Action $\rightarrow$ Observation};
un \emph{Auditor de Hechos}, que puede rechazar la respuesta y devolver el flujo con el motivo; y un \emph{Redactor}, que entrega la respuesta aprobada.
El orquestador es Mistral-7B-Instruct-v0.3 cuantizado a 4 bits y ejecutado en una GPU Tesla T4.
Las herramientas son una base vectorial ChromaDB sobre WikiQA, búsqueda web con DuckDuckGo, una calculadora segura y la ejecución aislada de Python.
Sobre 50 consultas (20 de WikiQA, 15 de GSM8K y 15 de MBPP) comparamos el sistema con un Naive RAG de un solo paso.
Medimos \emph{Faithfulness} y \emph{Answer Relevance} con LLM-as-a-Judge y, en GSM8K y MBPP, la corrección de forma objetiva.
El multi-agente logra una exactitud de 0.50 frente a 0.30 y entrega tres veces menos respuestas incorrectas (7 frente a 21 de 30).
En MBPP, donde el Auditor verifica ejecutando el código, no entregó ningún código erróneo.
En cambio, obtiene menor fidelidad y relevancia promediadas, porque sus abstenciones puntúan cero y porque el juez de 7B casi no distingue una respuesta correcta de una incorrecta.
Concluimos que un guardrail es tan bueno como la verificación en la que se apoya, y que las métricas de un juez débil deben complementarse con una verificación objetiva.
\end{abstract}

\section{Introducción}

Los modelos de lenguaje de tamaño medio, como los de 7 a 8 mil millones de parámetros, pueden ejecutarse localmente en una sola GPU, pero alucinan con frecuencia: afirman datos que no están en su contexto, razonan mal en problemas de varios pasos y escriben código que no funciona.
Un sistema de recuperación aumentada de un solo paso (\emph{Naive RAG}) mitiga el primer problema solo parcialmente y no ataca los otros dos.

Este proyecto diseña e implementa un sistema multi-agente autónomo que combina razonamiento y acción (ReAct) con un mecanismo de control explícito.
El objetivo no es solo responder, sino \emph{no propagar} respuestas sin respaldo.
Para eso, un agente auditor revisa el trabajo del agente investigador y puede rechazarlo, devolviendo el flujo con el motivo del rechazo, o abstenerse tras varios intentos.

Los objetivos concretos son:
\begin{enumerate}
  \item Definir al menos tres agentes con roles diferenciados en un grafo de control cíclico (LangGraph).
  \item Integrar herramientas reales mediante la sintaxis \emph{Thought $\rightarrow$ Action $\rightarrow$ Observation}: una base vectorial local y búsqueda web, además de una calculadora y un ejecutor de código.
  \item Implementar guardrails que permitan al Auditor rechazar de forma condicional el reporte del Investigador, y ejecutar en paralelo las consultas independientes.
  \item Evaluar cuantitativamente el sistema sobre 50 consultas con \emph{Faithfulness} y \emph{Answer Relevance}, y compararlo con un Naive RAG.
\end{enumerate}

El resto del informe se organiza así:
\begin{itemize}
  \item La sección~\ref{sec:marco} resume los conceptos base.
  \item La sección~\ref{sec:datos} describe los datos.
  \item La sección~\ref{sec:arquitectura} presenta la arquitectura.
  \item La sección~\ref{sec:experimentos} detalla la configuración experimental y las métricas.
  \item Las secciones~\ref{sec:resultados} y~\ref{sec:analisis} presentan los resultados y el análisis de errores.
  \item La sección~\ref{sec:conclusiones} concluye.
\end{itemize}

\section{Marco conceptual}\label{sec:marco}

\paragraph{ReAct.}
Yao et al.~\cite{react} proponen intercalar trazas de razonamiento (\emph{Thought}) con acciones sobre el entorno (\emph{Action}), cuyo resultado (\emph{Observation}) se agrega al contexto antes del siguiente paso.
Así el modelo puede buscar información, calcular o ejecutar código en vez de depender solo de lo que memorizó.

\paragraph{Recuperación aumentada (RAG).}
Un sistema RAG recupera fragmentos relevantes de una base de conocimiento y los agrega al prompt.
En su forma más simple (\emph{Naive RAG}) hace una sola recuperación y una sola llamada al modelo, sin verificar la respuesta.

\paragraph{Sistemas multi-agente y guardrails.}
Dividir el trabajo entre agentes con roles distintos permite separar \emph{producir} una respuesta de \emph{verificarla}.
Un guardrail es una regla de control que bloquea o devuelve una salida que no cumple ciertas condiciones.
Puede ser determinista, como ejecutar pruebas unitarias, o basarse en otro modelo que juzga la salida.

\paragraph{Evaluación con RAGAS y LLM-as-a-Judge.}
RAGAS~\cite{ragas} define métricas sin referencia para sistemas RAG.
La \emph{Faithfulness} es la fracción de afirmaciones de la respuesta que se pueden inferir del contexto recuperado.
La \emph{Answer Relevance} mide qué tan directamente la respuesta atiende la pregunta.
Ambas se calculan pidiéndole a un LLM que actúe como juez, de modo que heredan sus sesgos y su ruido.

\section{Datos}\label{sec:datos}

Se usan tres conjuntos públicos de Hugging Face, uno por tipo de tarea (tabla~\ref{tab:datos}).

\begin{table}[ht]
\centering
\caption{Conjuntos de datos y uso en el proyecto.}\label{tab:datos}
\small
\begin{tabularx}{\linewidth}{@{}l l X c@{}}
\toprule
Dataset & Partición & Uso & Consultas \\
\midrule
WikiQA~\cite{wikiqa} & \code{validation} & Base de conocimiento en ChromaDB y preguntas abiertas & 20 \\
GSM8K~\cite{gsm8k} & \code{test} & Problemas matemáticos de varios pasos; referencia numérica tras \code{\#\#\#\#} & 15 \\
MBPP~\cite{mbpp} & \code{test} & Generación de funciones Python; verificación con \code{test\_list} & 15 \\
\midrule
\multicolumn{3}{@{}l}{Total} & 50 \\
\bottomrule
\end{tabularx}
\end{table}

\paragraph{Base de conocimiento.}
WikiQA asocia cada pregunta con varias frases candidatas de Wikipedia, marcadas como correctas (\code{label=1}) o incorrectas (\code{label=0}).
Se agrupan las filas por pregunta y se conservan las 126 preguntas de validación con al menos una frase correcta.
Se indexan \emph{todas} sus frases candidatas: 1119 frases únicas, también las incorrectas.
Indexar solo las correctas pondría la respuesta al alcance de cualquier recuperación y favorecería artificialmente al Naive RAG.
Con las incorrectas, la base contiene distractores temáticamente cercanos.
El índice es una colección persistente de ChromaDB con su función de embeddings por defecto.
Los identificadores salen de un hash del contenido, así que reindexar no duplica documentos.

\paragraph{Consultas de MBPP.}
Siguiendo a Austin et al.~\cite{mbpp}, la pregunta incluye los \code{assert} de \code{test\_list} para que el modelo conozca el nombre y la firma de la función.
Ambos métodos ven esas pruebas, pero solo el multi-agente puede ejecutarlas.
Su exactitud en MBPP mide, por lo tanto, si sabe verificar y corregir su código, no si generaliza a pruebas ocultas.

\section{Arquitectura del sistema}\label{sec:arquitectura}

\subsection{Grafo de control}

La figura~\ref{fig:grafo} muestra el grafo de LangGraph.
Su estado compartido contiene la pregunta, la evidencia recuperada, el reporte del Investigador, el veredicto del Auditor, la retroalimentación, el número de rondas y la respuesta final.

\begin{figure}[ht]
\centering
\begin{tikzpicture}[
  node distance=0.9cm and 1.1cm,
  >={Stealth[length=2.2mm]},
  every node/.style={font=\small},
  agent/.style={draw, rounded corners, fill=multi!12, minimum height=1.1cm, minimum width=2.0cm, align=center},
  tool/.style={draw, fill=gray!10, minimum height=0.8cm, minimum width=1.6cm, align=center},
  term/.style={draw, ellipse, minimum height=0.7cm, align=center},
  decision/.style={draw, diamond, aspect=1.6, fill=naive!15, inner sep=1pt, align=center},
]
  \node[term] (start) {inicio};
  \node[tool, minimum width=2.0cm, above right=0.25cm and 0.55cm of start] (kb) {\code{kb}\\ ChromaDB};
  \node[tool, minimum width=2.0cm, below right=0.25cm and 0.55cm of start] (web) {\code{web}\\ DuckDuckGo};
  % El Investigador se ubica a partir del borde derecho de kb y web, así la unión no las toca.
  \coordinate (merge) at ([xshift=4mm]web.east |- start);
  \node[agent, right=4mm of merge] (inv) {Investigador\\ \footnotesize ReAct ($\leq 6$ pasos)};
  \node[decision, right=0.8cm of inv] (aud) {Auditor\\de Hechos};
  \node[agent, right=1.1cm of aud] (red) {Redactor};
  \node[term, right=0.4cm of red] (end) {fin};
  \node[tool, below=1.3cm of aud] (fb) {feedback\\ \footnotesize motivo + respuesta};
  \node[term, above=1.1cm of aud] (abs) {abstención};

  \draw[->] (start) |- (kb);
  \draw[->] (start) |- (web);
  \draw (kb.east) -| (merge);
  \draw (web.east) -| (merge);
  \draw[->] (merge) -- (inv.west);
  \draw[->] (inv) -- (aud);
  \draw[->] (aud) -- node[above, pos=0.42, font=\footnotesize] {aprobado} (red);
  \draw[->] (red) -- (end);
  \draw[->] (aud) -- node[right, font=\footnotesize] {rechazo} (fb);
  \draw[->] (fb) -| (inv);
  \draw[->] (aud) -- node[right, font=\footnotesize, align=left] {$\geq$ \code{max\_rounds}\\rechazos} (abs);
  \node[font=\footnotesize, gray, below=0.05cm of web] {en paralelo (fan-out)};
\end{tikzpicture}
\caption{Grafo cíclico de control. \code{kb} y \code{web} se consultan en paralelo antes del bucle ReAct (solo en WikiQA). Si el Auditor rechaza, el motivo y la respuesta rechazada vuelven al Investigador. Tras \code{max\_rounds}~$=3$ rechazos el sistema se abstiene.}
\label{fig:grafo}
\end{figure}

\subsection{Agentes}

\paragraph{Investigador.}
Ejecuta un bucle ReAct de hasta seis pasos.
En cada paso el LLM escribe un \emph{Thought} y una sola \emph{Action} con la forma \code{herramienta[argumento]}.
Solo el sistema escribe la \emph{Observation}: si el modelo inventa una, se descarta todo lo que viene después.
El bucle termina cuando el modelo escribe \code{Final:}.
El prompt incluye un ejemplo resuelto propio de cada tarea, la retroalimentación del Auditor y las observaciones acumuladas.

\paragraph{Auditor de Hechos.}
Es el guardrail del sistema y aplica las reglas en este orden:
\begin{enumerate}
  \item Rechaza si no hubo respuesta final dentro del límite de pasos.
  \item En código, ejecuta la respuesta junto con los \code{assert} del enunciado y la aprueba solo si pasa; esta verificación no consulta al LLM.
  \item Rechaza si la respuesta no tiene evidencia, o si incluye código que no se ejecutó tal cual.
  \item En los demás casos, le pide al LLM que responda \code{APROBADO} o \code{RECHAZADO: motivo}, con una regla específica de la tarea. Por ejemplo, en GSM8K le pide verificar que el número final responda exactamente lo que se pregunta.
\end{enumerate}

\paragraph{Redactor.}
Reformula la respuesta aprobada en una a tres oraciones, sin agregar datos ni fuentes.
En matemáticas termina con \code{Respuesta final: <número>}.
El código ya verificado se entrega sin cambios: reescribirlo con un LLM solo podría estropearlo.

\subsection{Herramientas por tarea}

Se construye un grafo por tipo de tarea con \code{build\_flow(..., task=...)}.
Cada grafo tiene sus herramientas (tabla~\ref{tab:herramientas}), su ejemplo resuelto y sus reglas de auditoría y redacción.
En GSM8K y MBPP no se ofrece la web: no aporta y podría devolver la solución publicada del propio benchmark.

\begin{table}[ht]
\centering
\caption{Herramientas disponibles para el Investigador en cada tarea.}\label{tab:herramientas}
\small
\begin{tabularx}{\linewidth}{@{}l l X@{}}
\toprule
Tarea & Herramienta & Implementación \\
\midrule
\multirow{2}{*}{WikiQA} & \code{kb[...]} & Búsqueda semántica en ChromaDB (3 fragmentos) \\
 & \code{web[...]} & DuckDuckGo Search \\
GSM8K & \code{calc[...]} & Evaluador aritmético sobre el árbol sintáctico de Python, sin \code{eval}, con límite de exponente \\
MBPP & \code{python[...]} & Proceso aislado (\code{python -I}, entorno vacío, directorio temporal) con límites de tiempo, CPU y memoria \\
\bottomrule
\end{tabularx}
\end{table}

\subsection{Concurrencia}

En WikiQA, \code{kb} y \code{web} son consultas independientes.
El grafo las lanza en paralelo con un \emph{fan-out} de LangGraph.
Las dos ramas escriben en el mismo campo del estado, que usa un reductor de concatenación, y el Investigador espera a ambas antes de empezar.
Si una rama falla, devuelve una lista vacía y no detiene a la otra.
Una prueba con una barrera de sincronización comprueba que ambas ramas se ejecutan a la vez.

\subsection{Ajustes por los límites del modelo de 7B}

Las primeras pruebas en la T4 mostraron comportamientos que se corrigieron en el diseño:
\begin{itemize}
  \item \textbf{Acciones repetidas en bucle.} Se agregaron una regla anti-repetición (repetir una \emph{Action} devuelve un aviso en vez de ejecutarla), un ejemplo resuelto y un último paso que exige \code{Final}.
  \item \textbf{Código transcrito con errores.} Si el modelo dice que su código pasó las pruebas pero no lo copia en \code{Final}, se entrega el último código que realmente ejecutó.
  \item \textbf{Relleno en la redacción.} El Redactor inventaba explicaciones que bajaban la fidelidad, así que su prompt prohíbe agregar contenido.
  \item \textbf{Idioma.} La base y la web están en inglés, y las búsquedas en español recuperaban fragmentos irrelevantes. El Investigador busca en el idioma de la pregunta.
\end{itemize}

\subsection{Arquitectura de software}

El código sigue una arquitectura hexagonal.
El núcleo (\code{domain} y \code{application}) solo depende de puertos (\code{LLMPort}, \code{KnowledgeBasePort}, \code{WebSearchPort}, \code{CodeRunnerPort}, entre otros), y los adaptadores implementan cada puerto.
Cambiar Mistral por otro LLM, o ChromaDB por FAISS, es reemplazar un adaptador.
Esto permite además probar todo el flujo sin GPU ni red con adaptadores falsos: el repositorio tiene 59 pruebas automáticas.

\section{Configuración experimental}\label{sec:experimentos}

\subsection{Modelo y hardware}

El orquestador de los tres agentes es \code{mistralai/Mistral-7B-Instruct-v0.3}~\cite{mistral}, cuantizado a 4 bits (NF4~\cite{qlora}, con cómputo en \code{float16}) mediante \code{bitsandbytes}.
Se ejecuta con \code{transformers} en una GPU Tesla T4 (16\,GB) de Google Colab.
La decodificación es voraz (\code{do\_sample=False}), con hasta 512 tokens nuevos, y se detiene en la cadena \code{Observation:}.
La corrida completa, de 100 registros (50 consultas con dos métodos), tomó unos 96 minutos.

\subsection{Sistemas comparados}

\begin{itemize}
  \item \textbf{Multi-agente.} Es el grafo de la sección~\ref{sec:arquitectura}, con \code{max\_rounds}~$=3$ y hasta 6 pasos ReAct por ronda.
  \item \textbf{Naive RAG.} Recupera 3 fragmentos de la misma base ChromaDB y hace una sola llamada al mismo modelo, sin herramientas ni auditoría.
  En GSM8K y MBPP la base de WikiQA no aporta contexto útil, así que ahí la comparación es contra un modelo que razona sin herramientas.
\end{itemize}

Ambos sistemas reciben exactamente las mismas 50 preguntas y los evalúa el mismo juez.

\subsection{Métricas}

Sea $q$ una pregunta, $a$ la respuesta y $C$ el contexto del método: los fragmentos recuperados y, en el multi-agente, también las salidas de \code{calc} y \code{python}.

\paragraph{Faithfulness.}
El juez descompone $a$ en a lo sumo 8 afirmaciones atómicas $s_1,\dots,s_n$ y marca cada una como inferible (\code{SI}) o no (\code{NO}) a partir de $C$ y de los datos de $q$:
\begin{equation}
  \text{Faithfulness}(a, C) = \frac{\left|\{\,s_i : s_i \text{ se infiere de } C \cup q\,\}\right|}{n}.
\end{equation}
Los datos del enunciado cuentan como información dada porque en GSM8K y MBPP la respuesta los repite con razón.

\paragraph{Answer Relevance.}
El juez asigna un entero $r \in \{0,\dots,10\}$ según qué tan directamente $a$ responde $q$, sin juzgar si es verdadera, y se normaliza como $r/10$.
Es una simplificación de RAGAS, que genera preguntas a partir de la respuesta y mide su similitud con la original.

\paragraph{Verificación objetiva.}
En GSM8K se extrae el número que sigue a \code{Respuesta final:} (o, si no existe esa línea, el último número del texto) y se compara con la referencia.
En MBPP se ejecuta el código de la respuesta con sus \code{assert}.
Sobre estas 30 consultas se reportan tres métricas:
\begin{align}
  \text{exactitud} &= \frac{\#\,\text{correctas}}{\#\,\text{consultas}}, \\
  \text{exactitud al responder} &= \frac{\#\,\text{correctas}}{\#\,\text{respondidas}}, \\
  \text{incorrectas entregadas} &= \frac{\#\,\text{respondidas e incorrectas}}{\#\,\text{consultas}}.
\end{align}
La última mide cuántas alucinaciones llegan al usuario: una abstención no cuenta como incorrecta.

\paragraph{Abstenciones.}
Si el Auditor nunca aprueba, el multi-agente no responde.
Una abstención puntúa 0 en \emph{Faithfulness} y \emph{Answer Relevance}, y cuenta como no correcta en la exactitud.
Para separar ``qué tan fiel es cuando responde'' de ``cuántas veces se abstiene'', se reporta además la fidelidad promediada solo sobre las respuestas entregadas.

\paragraph{Juez.}
El juez es el mismo Mistral-7B que orquesta.
Así la evaluación cabe en la misma GPU, pero el juez es ruidoso y podría favorecer sus propias respuestas.
Si el juez no devuelve el formato pedido se reintenta una vez; si vuelve a fallar, la métrica queda vacía y se cuenta aparte.

\paragraph{Prueba estadística.}
Para la exactitud se usa la prueba exacta de McNemar sobre los pares discordantes: las consultas que acierta un solo método.
Es la prueba adecuada para dos clasificadores evaluados sobre las mismas muestras.

\section{Resultados}\label{sec:resultados}

La tabla~\ref{tab:global} resume la corrida sobre las 50 consultas.
La tabla~\ref{tab:dataset} la desglosa por dataset.

\begin{table}[ht]
\centering
\caption{Resultados globales. Las métricas del juez promedian las 50 consultas; la verificación objetiva, las 30 de GSM8K y MBPP. En negrita, el mejor valor; $\downarrow$ indica que menor es mejor.}\label{tab:global}
\small
\begin{tabular}{@{}lcc@{}}
\toprule
Métrica & Multi-agente & Naive RAG \\
\midrule
\multicolumn{3}{@{}l}{\emph{Métricas del juez (50 consultas)}} \\
\quad Respondidas & 0.84 & \textbf{1.00} \\
\quad Faithfulness & 0.580 & \textbf{0.663} \\
\quad Faithfulness (solo respondidas) & \textbf{0.691} & 0.663 \\
\quad Answer Relevance & 0.792 & \textbf{0.911} \\
\midrule
\multicolumn{3}{@{}l}{\emph{Verificación objetiva (30 consultas)}} \\
\quad Exactitud & \textbf{0.50} (15/30) & 0.30 (9/30) \\
\quad Exactitud al responder & \textbf{0.68} (15/22) & 0.30 (9/30) \\
\quad Incorrectas entregadas $\downarrow$ & \textbf{0.23} (7/30) & 0.70 (21/30) \\
\midrule
Rondas promedio & 1.40 & 1.00 \\
Segundos por consulta & 88.9 & \textbf{26.0} \\
\bottomrule
\end{tabular}
\end{table}

\begin{table}[ht]
\centering
\caption{Resultados por dataset. WikiQA no tiene verificación objetiva.}\label{tab:dataset}
\small
\setlength{\tabcolsep}{4pt}
\begin{tabular}{@{}llcccccc@{}}
\toprule
Dataset & Método & Resp. & Faith. & Faith. (resp.) & Relev. & Exactitud & Incorr. entr. $\downarrow$ \\
\midrule
\multirow{2}{*}{WikiQA}
 & Multi-agente & 1.00 & 0.605 & 0.605 & 0.88 & --- & --- \\
 & Naive RAG    & 1.00 & 0.649 & 0.649 & 0.89 & --- & --- \\
\midrule
\multirow{2}{*}{GSM8K}
 & Multi-agente & 0.80 & 0.679 & 0.849 & 0.80 & \textbf{0.33} & \textbf{0.47} \\
 & Naive RAG    & 1.00 & 0.827 & 0.827 & 0.91 & 0.20 & 0.80 \\
\midrule
\multirow{2}{*}{MBPP}
 & Multi-agente & 0.67 & 0.447 & 0.671 & 0.67 & \textbf{0.67} & \textbf{0.00} \\
 & Naive RAG    & 1.00 & 0.519 & 0.519 & 0.95 & 0.40 & 0.60 \\
\bottomrule
\end{tabular}
\end{table}

\begin{figure}[ht]
\centering
\begin{tikzpicture}
\begin{axis}[
  ybar, bar width=11pt, width=0.48\linewidth, height=5.2cm,
  ymin=0, ymax=1.05, ylabel={Exactitud},
  symbolic x coords={GSM8K,MBPP,Total}, xtick=data,
  nodes near coords, nodes near coords style={font=\tiny},
  every node near coord/.append style={/pgf/number format/fixed, /pgf/number format/precision=2, /pgf/number format/fixed zerofill},
  legend style={at={(0.5,1.03)}, anchor=south, legend columns=2, font=\footnotesize, draw=none, /tikz/every even column/.append style={column sep=8pt}},
  tick label style={font=\footnotesize}, label style={font=\footnotesize},
  enlarge x limits=0.25,
]
\addplot[fill=multi, draw=multi] coordinates {(GSM8K,0.333) (MBPP,0.667) (Total,0.50)};
\addplot[fill=naive, draw=naive] coordinates {(GSM8K,0.20) (MBPP,0.40) (Total,0.30)};
\legend{Multi-agente, Naive RAG}
\end{axis}
\end{tikzpicture}\hfill
\begin{tikzpicture}
\begin{axis}[
  ybar, bar width=11pt, width=0.48\linewidth, height=5.2cm,
  ymin=0, ymax=1.05, ylabel={Incorrectas entregadas ($\downarrow$)},
  symbolic x coords={GSM8K,MBPP,Total}, xtick=data,
  nodes near coords, nodes near coords style={font=\tiny},
  every node near coord/.append style={/pgf/number format/fixed, /pgf/number format/precision=2, /pgf/number format/fixed zerofill},
  legend style={at={(0.5,1.03)}, anchor=south, legend columns=2, font=\footnotesize, draw=none, /tikz/every even column/.append style={column sep=8pt}},
  tick label style={font=\footnotesize}, label style={font=\footnotesize},
  enlarge x limits=0.25,
]
\addplot[fill=multi, draw=multi] coordinates {(GSM8K,0.467) (MBPP,0.0) (Total,0.233)};
\addplot[fill=naive, draw=naive] coordinates {(GSM8K,0.80) (MBPP,0.60) (Total,0.70)};
\legend{Multi-agente, Naive RAG}
\end{axis}
\end{tikzpicture}
\caption{Verificación objetiva por dataset. Izquierda: fracción de consultas resueltas correctamente. Derecha: fracción de consultas con una respuesta entregada y errónea.}
\label{fig:objetiva}
\end{figure}

\paragraph{Lectura de los resultados.}
\begin{itemize}
  \item \textbf{Verificación objetiva.} El multi-agente supera al Naive RAG en todos los indicadores (figura~\ref{fig:objetiva}).
  Resuelve 15 de 30 consultas frente a 9, y entrega 7 respuestas incorrectas frente a 21.
  En MBPP no entregó ningún código erróneo.
  \item \textbf{Métricas del juez.} El Naive RAG obtiene mayor \emph{Faithfulness} y \emph{Answer Relevance} promediadas sobre las 50 consultas.
  La diferencia se explica sobre todo por las 8 abstenciones del multi-agente (3 en GSM8K y 5 en MBPP), que puntúan 0.
  Si se promedian solo las respuestas entregadas, el multi-agente es más fiel (0.691 frente a 0.663).
  \item \textbf{WikiQA.} Los dos métodos casi empatan: fidelidad de 0.605 frente a 0.649 y relevancia de 0.88 frente a 0.89.
  \item \textbf{Costo.} El multi-agente es unas 3.4 veces más lento, por las varias llamadas al LLM de cada consulta.
\end{itemize}

\section{Análisis de errores}\label{sec:analisis}

Las cifras de esta sección se reproducen con \code{scripts/analisis\_errores.py}, a partir del archivo con los 100 registros de la corrida.

\subsection{Comparación consulta por consulta}

\begin{table}[ht]
\centering
\caption{Resultado de cada método en las 30 consultas con verificación objetiva.}\label{tab:pares}
\small
\begin{tabular}{@{}llccc@{}}
\toprule
Multi-agente & Naive RAG & GSM8K & MBPP & Total \\
\midrule
correcta   & correcta   & 1 & 6 & 7 \\
correcta   & incorrecta & 4 & 4 & \textbf{8} \\
incorrecta & correcta   & 1 & 0 & \textbf{1} \\
incorrecta & incorrecta & 6 & 0 & 6 \\
abstención & correcta   & 1 & 0 & 1 \\
abstención & incorrecta & 2 & 5 & 7 \\
\bottomrule
\end{tabular}
\end{table}

La tabla~\ref{tab:pares} muestra que el multi-agente acierta 8 consultas que el Naive RAG falla.
El Naive RAG acierta solo 2 que el multi-agente no acierta: una por error y otra por abstención.
Con la prueba exacta de McNemar sobre esos 10 pares discordantes se obtiene $p = 0.109$.
La dirección del efecto es consistente, pero 30 consultas no alcanzan para descartar el azar con un nivel de 0.05.

Las abstenciones cumplen su propósito: en 7 de las 8, el Naive RAG entregó una respuesta incorrecta.
El guardrail cambió respuestas equivocadas por abstenciones, que es exactamente lo que se busca para no propagar alucinaciones.

\subsection{El Auditor es tan bueno como su verificación}

\begin{table}[ht]
\centering
\caption{Qué pasó con las consultas que el Auditor rechazó al menos una vez.}\label{tab:auditor}
\small
\begin{tabular}{@{}lcccc@{}}
\toprule
Dataset & Con rechazo & Corregidas & Abstenciones & Aprobadas pero incorrectas \\
\midrule
WikiQA & 1 de 20 & 1 & 0 & --- \\
GSM8K  & 3 de 15 & 0 & 3 & 7 \\
MBPP   & 7 de 15 & 2 & 5 & \textbf{0} \\
\bottomrule
\end{tabular}
\end{table}

La tabla~\ref{tab:auditor} contrasta los dos tipos de guardrail:
\begin{itemize}
  \item \textbf{MBPP (verificación determinista).} El Auditor ejecuta el código con sus pruebas y nunca aprobó un código erróneo.
  En dos consultas el error de ejecución devuelto al Investigador le permitió corregir el código, en la segunda y en la tercera ronda.
  El Naive RAG falló ambas.
  \item \textbf{GSM8K (verificación con el LLM).} El Auditor solo puede leer el razonamiento, y el mismo modelo que se equivocó tiende a aprobar su propio error.
  De las 7 respuestas aprobadas que el verificador marcó como incorrectas, 6 son errores de razonamiento y no de aritmética, porque \code{calc} calcula bien lo que se le pide.
  Por ejemplo, interpreta ``el valor aumentó un 150\,\%'' como multiplicar por 1.5 en vez de por 2.5, o plantea bien la suma $160+80+20$ y luego responde 240 sin pasarla por \code{calc}.
\end{itemize}

\subsection{El juez frente a la verificación objetiva}

Sobre las 52 respuestas entregadas en GSM8K y MBPP por ambos métodos, el juez asigna casi el mismo puntaje a las correctas que a las incorrectas (tabla~\ref{tab:juez}).

\begin{table}[ht]
\centering
\caption{Puntaje promedio del juez según la verificación objetiva de la respuesta.}\label{tab:juez}
\small
\begin{tabular}{@{}lccc@{}}
\toprule
Respuestas & $n$ & Answer Relevance & Faithfulness \\
\midrule
Correctas   & 24 & 1.000 & 0.735 \\
Incorrectas & 28 & 0.927 & 0.694 \\
\bottomrule
\end{tabular}
\end{table}

En \emph{Answer Relevance} esto es esperable, porque la métrica no juzga la veracidad.
En \emph{Faithfulness} el juez también falla: un código correcto que pasó todas sus pruebas recibió fidelidad 0.0, y dos códigos incorrectos del Naive RAG recibieron 1.0.
Con un juez de 7B estas métricas separan mal a los dos métodos.
Por eso la verificación objetiva es la evidencia más sólida de la comparación.

\subsection{Otros hallazgos}

\begin{itemize}
  \item \textbf{Falso negativo del verificador.} En una consulta de GSM8K el multi-agente respondió correctamente ``\$64'', pero el Redactor no escribió la línea \code{Respuesta final:} y el verificador tomó el último número del texto.
  Se revisó a mano el número extraído en las 30 respuestas de GSM8K y es el único caso.
  Los resultados reportados no se corrigieron; con la corrección, la exactitud del multi-agente en GSM8K sería 0.40 en vez de 0.33.
  \item \textbf{Fidelidad baja en respuestas cortas de WikiQA.} Respuestas correctas como ``Holly Ridge, North Carolina is located in Onslow County.'' recibieron fidelidad de 0.12.
  Las respuestas cortas del multi-agente promedian 0.48, frente a 0.71 de las largas.
  La hipótesis más probable es que el juez rinde peor con el contexto más largo del multi-agente (7 fragmentos en promedio frente a 3), pero no se comprobó.
  \item \textbf{Idioma.} El Redactor respondió en español a 7 de las 20 preguntas en inglés de WikiQA, probablemente porque su prompt está en español.
  Este efecto no explica la fidelidad baja: esas respuestas puntúan más que las respuestas en inglés.
\end{itemize}

\section{Discusión y conclusiones}\label{sec:conclusiones}

\paragraph{Cumplimiento de los objetivos.}
El sistema implementa tres agentes con roles interdependientes en un grafo cíclico de LangGraph.
El Investigador sigue el ciclo \emph{Thought $\rightarrow$ Action $\rightarrow$ Observation} con una base vectorial local (ChromaDB sobre WikiQA), búsqueda web (DuckDuckGo), una calculadora y un ejecutor de código.
El Auditor rechaza de forma condicional y devuelve el flujo con el motivo, y las consultas independientes se ejecutan en paralelo.
Todo corre con un modelo de 7B cuantizado a 4 bits en una T4.

\paragraph{¿Supera el multi-agente al Naive RAG?}
Donde la corrección se puede verificar, sí:
\begin{itemize}
  \item Logra 20 puntos más de exactitud (0.50 frente a 0.30).
  \item Entrega tres veces menos respuestas incorrectas (7 frente a 21 de 30).
  \item En MBPP no entregó ningún código erróneo.
\end{itemize}
En las métricas del juez promediadas sobre todas las consultas, no.
Las abstenciones puntúan cero, y un juez de 7B apenas distingue una respuesta correcta de una incorrecta.
Si se consideran solo las respuestas entregadas, el multi-agente es más fiel.

\paragraph{Lecciones.}
\begin{enumerate}
  \item \textbf{Un guardrail vale lo que vale su verificación.} Ejecutar pruebas es un guardrail fuerte, y pedirle al mismo modelo que revise su razonamiento es uno débil.
  Conviene usar verificaciones deterministas siempre que la tarea lo permita.
  \item \textbf{Abstenerse es una salida válida.} En un contexto de auditoría, no responder es preferible a responder mal.
  Las métricas que promedian incluyendo las abstenciones como cero castigan justamente ese comportamiento, así que deben acompañarse de métricas como las respuestas incorrectas entregadas.
  \item \textbf{El juez condiciona la conclusión.} Con un juez débil y el mismo que responde, \emph{Faithfulness} y \emph{Answer Relevance} deben leerse con cautela.
  \item \textbf{Un modelo de 7B necesita andamiaje.} La regla anti-repetición, los ejemplos resueltos por tarea, entregar el código realmente ejecutado y un Redactor que no reescribe fueron necesarios para que el flujo funcionara.
\end{enumerate}

\paragraph{Limitaciones.}
\begin{itemize}
  \item Es una sola corrida de 50 consultas, y la ventaja en exactitud no alcanza significancia al 5\,\% ($p = 0.109$).
  \item El juez es el mismo modelo que responde.
  \item Los registros no guardan el texto de la evidencia ni el motivo de cada rechazo, lo que impide volver a juzgar sin repetir la generación.
  \item El multi-agente es unas 3.4 veces más lento.
  \item El ejecutor de código limita tiempo, CPU y memoria, pero no aísla la red ni el sistema de archivos. Es aceptable en una máquina virtual desechable como la de Colab.
\end{itemize}

\paragraph{Trabajo futuro.}
\begin{enumerate}
  \item Usar un juez independiente y más fuerte.
  \item Guardar la evidencia en cada registro para volver a juzgar sin volver a generar.
  \item Hacer una ablación del Auditor (\code{max\_rounds}~$=1$) y de la recuperación en paralelo, para medir el aporte de cada componente.
  \item Ampliar a unas 100 consultas con verificación objetiva para obtener significancia estadística.
  \item Agregar en GSM8K una verificación determinista, por ejemplo pedir el planteamiento como una expresión que \code{calc} pueda reevaluar.
\end{enumerate}

\begin{thebibliography}{9}

\bibitem{react}
S.~Yao, J.~Zhao, D.~Yu, N.~Du, I.~Shafran, K.~Narasimhan y Y.~Cao.
\newblock ReAct: Synergizing reasoning and acting in language models.
\newblock En \emph{International Conference on Learning Representations (ICLR)}, 2023.

\bibitem{wikiqa}
Y.~Yang, W.~Yih y C.~Meek.
\newblock WikiQA: A challenge dataset for open-domain question answering.
\newblock En \emph{Proceedings of EMNLP}, 2015.

\bibitem{gsm8k}
K.~Cobbe, V.~Kosaraju, M.~Bavarian, M.~Chen, H.~Jun, L.~Kaiser, M.~Plappert, J.~Tworek, J.~Hilton, R.~Nakano, C.~Hesse y J.~Schulman.
\newblock Training verifiers to solve math word problems.
\newblock \emph{arXiv:2110.14168}, 2021.

\bibitem{mbpp}
J.~Austin, A.~Odena, M.~Nye, M.~Bosma, H.~Michalewski, D.~Dohan, E.~Jiang, C.~Cai, M.~Terry, Q.~Le y C.~Sutton.
\newblock Program synthesis with large language models.
\newblock \emph{arXiv:2108.07732}, 2021.

\bibitem{ragas}
S.~Es, J.~James, L.~Espinosa-Anke y S.~Schockaert.
\newblock RAGAS: Automated evaluation of retrieval augmented generation.
\newblock En \emph{Proceedings of EACL: System Demonstrations}, 2024.

\bibitem{mistral}
A.~Q.~Jiang et~al.
\newblock Mistral 7B.
\newblock \emph{arXiv:2310.06825}, 2023.

\bibitem{qlora}
T.~Dettmers, A.~Pagnoni, A.~Holtzman y L.~Zettlemoyer.
\newblock QLoRA: Efficient finetuning of quantized LLMs.
\newblock En \emph{Advances in Neural Information Processing Systems (NeurIPS)}, 2023.

\end{thebibliography}


\end{document}
